%% Correction audit de veracite - 2026-09-20 : citations de rapport de jury sans
%% support textuel dans le corpus (rapports 2018-2025) remplacees par des verbatims
%% verifies ou supprimees. Detail ci-dessous.
%% l.336 (ex) : citation attribuee a un "Rapport jury 2026" INEXISTANT (le corpus
%%   s'arrete a 2025) et INTROUVABLE -> remplacee par RJ2023 l.459 (verbatim, cite 2023).
%% l.337 (ex) : citation "jury 2024" sur le filtre passe-bas RL des enroulements :
%%   INTROUVABLE comme commentaire de jury (le fait figure au corrige technique
%%   RJ2024 l.239) -> requalifiee explicitement en extrait du corrige officiel.
\expandafter\ifx\csname CAPETINCLUDEMODE\endcsname\relax
\documentclass[11pt,a4paper]{article}
% preamble.tex — packages et commandes communes a toutes les fiches CAPET SII IE
\input{glyphtounicode}
\pdfgentounicode=1
\usepackage{cmap}
\usepackage[T1]{fontenc}
\usepackage[utf8]{inputenc}
\usepackage[french]{babel}
% babel[french] active a \begin{document} des raccourcis de ponctuation (:;!?)
% qui cassent la syntaxe to[...] de circuitikz (erreur "+ or - expected") ;
% on les desactive, ce qui ne change que l'espacement fin avant ces signes.
\AtBeginDocument{\shorthandoff{:;!?}}
\usepackage[a4paper, top=1.6cm, bottom=1.6cm, left=1.4cm, right=1.4cm]{geometry}
\setlength{\headheight}{14.5pt}
\usepackage{amsmath, amssymb, amsfonts}
\usepackage{siunitx}
\sisetup{locale=FR, output-decimal-marker={,}, inter-unit-product=\ensuremath{{}\cdot{}},
         group-minimum-digits=4}
% Unites non fournies par siunitx, declarees ici une fois pour toutes (fiches + exercices).
% Les gardes \ifdefined evitent tout conflit avec une redeclaration locale restee dans un
% fichier : ne pas retirer les gardes, ne pas redeclarer ces unites dans les fichiers.
\DeclareSIUnit{\jour}{j}
\ifdefined\var\else\DeclareSIUnit{\var}{var}\fi   % puissance reactive
\ifdefined\vad\else\DeclareSIUnit{\vad}{VAD}\fi   % puissance deformante
\ifdefined\tr\else\DeclareSIUnit{\tr}{tr}\fi      % tour (tr/min)
\usepackage{array, booktabs, multirow, tabularx}
\usepackage{xcolor}
\usepackage[most, breakable]{tcolorbox}
\usepackage{multicol}
\usepackage{enumitem}
\usepackage{mathtools}
\usepackage{lmodern}
% Refonte visuelle 2026 (direction "moderne coloré") : Lato pour le texte,
% maths sans-serif assorties via newtxsf (charge apres amsmath/amssymb).
\usepackage[default]{lato}
\usepackage{newtxsf}
\usepackage{microtype}
\usepackage{fancyhdr}
\usepackage{tikz}
\usetikzlibrary{positioning,arrows.meta}
\usepackage{circuitikz}
\usepackage{lastpage}

% ── Couleurs par domaine (palette modernisée 2026) ────────────────────
\definecolor{colElec}{RGB}{29, 78, 216}        % Électrotechnique — bleu vif
\definecolor{colElecLight}{RGB}{219, 234, 254} % fond bleu clair
\definecolor{colAuto}{RGB}{21, 128, 61}        % Automatique — vert
\definecolor{colMeca}{RGB}{194, 65, 12}        % Mécanique — orange
\definecolor{colTherm}{RGB}{185, 28, 28}       % Thermique — rouge
\definecolor{colQT}{RGB}{232, 244, 233}        % (hérité, plus utilisé par questiontype)
\definecolor{colJury}{RGB}{255, 248, 225}      % fond jury (crème)
\definecolor{colJuryFrame}{RGB}{202, 138, 4}   % bordure/bandeau jury (ambre)
\definecolor{colGrayLight}{RGB}{245, 245, 247} % fond gris neutre
\definecolor{colSeparator}{RGB}{180, 180, 180} % gris séparateur
\definecolor{colExo}{RGB}{8, 145, 178}          % Banque d'exercices — cyan (distinct des fiches)
% ── Encre et gris de la refonte ───────────────────────────────────────
\definecolor{colInk}{RGB}{23, 32, 42}          % texte courant (quasi-noir)
\definecolor{colSlate}{RGB}{83, 98, 118}       % gris ardoise (méta, en-têtes)
\definecolor{colRule}{RGB}{203, 213, 225}      % filets clairs
\definecolor{colJuryBg}{RGB}{255, 248, 225}    % fond crème jury
\definecolor{colPresence}{RGB}{83, 98, 118}    % bloc 5 (présence dans les sujets)
\AtBeginDocument{\color{colInk}}

% ── Style encadré exercice + corrigé (banque d'exercices) ─────────────
\tcbset{
  exoBox/.style={
    breakable,
    enhanced,
    colback=colExo!8!white,
    colframe=colExo,
    colbacktitle=colExo,
    coltitle=white,
    fonttitle=\bfseries\small,
    left=4pt, right=4pt, top=3pt, bottom=3pt,
    before skip=5pt, after skip=5pt,
  }
}

% ── Bascule enonce / corrige (banque d'exercices) ─────────────────────
% \CAPETCORRIGE est defini en ligne de commande (pdflatex "\def\CAPETCORRIGE{1}\input{...}")
% pour compiler la variante corriges d'un meme fichier source ; par defaut (non defini),
% la compilation produit la variante enonces. \CAPETSHOW{enonce}{corrige} choisit la branche.
%
% TROIS MODES (refonte 2026-09-20, manuel de revision) :
%   - aucun drapeau   -> imprime l'enonce            (variante enonces, INCHANGE)
%   - \CAPETCORRIGE   -> imprime le corrige          (variante corriges, INCHANGE)
%   - \CAPETMANUEL    -> imprime l'enonce ET met le corrige de cote pour l'annexe
%
% Le mode manuel stocke chaque corrige dans une macro (\csgdef), sans ecriture
% fichier ni re-tokenisation : le contenu (circuitikz, accents, %) traverse intact.
% L'annexe les rejoue ensuite dans l'ordre via \CAPETANNEXE.
\usepackage{etoolbox}

\newcounter{CAPETexo}          % numero continu d'exercice dans le manuel
\setcounter{CAPETexo}{0}
\newcounter{CAPETannexe}       % compteur de relecture en annexe

% Titre courant de la notion en cours (pose par \CAPETNOTION), sert d'intertitre annexe
\newcommand{\CAPETnotioncode}{}
\newcommand{\CAPETnotiontitre}{}

% \CAPETNOTION{code}{titre} : ouvre une notion dans le manuel (memorise le libelle,
% et marque le rang du premier exercice de la notion pour l'intertitre d'annexe).
\newcommand{\CAPETNOTION}[2]{%
  \renewcommand{\CAPETnotioncode}{#1}%
  \renewcommand{\CAPETnotiontitre}{#2}%
  \csgdef{CAPETnotionstart\the\numexpr\value{CAPETexo}+1\relax}{#1 -- #2}%
}

% ── \degree robuste ───────────────────────────────────────────────────
% Certains fichiers d'exercices utilisent \degree. Ni le preambule ni siunitx
% ne le definissent, et \providecommand local ne suffit pas : sous hyperref,
% \degree existe deja comme commande factice (\ERROR), donc \providecommand
% s'abstient et l'appel echoue a la compilation du manuel (erreur invisible
% en compilation isolee). On le (re)definit donc inconditionnellement ici,
% apres le chargement des packages, pour tout le corpus.
\AtBeginDocument{\providecommand{\degree}{}\renewcommand{\degree}{\ensuremath{^\circ}}}

% Libelle decoratif : les fichiers d'exercices utilisent aussi \CAPETSHOW dans
% leur \fancyhead et leur bloc de titre, pour afficher "(enonces)" ou "(corriges)".
% Ces appels-la ne sont PAS des exercices : en mode manuel ils ne doivent ni
% consommer un numero, ni alimenter l'annexe. On les reconnait a leur premier
% argument, strictement egal au libelle ci-dessous.
\newcommand{\CAPETLIBELLEENONCES}{(\'enonc\'es)}

\ifdefined\CAPETMANUEL
  % --- Mode manuel : enonce en place, corrige differe vers l'annexe ---
  \newcommand{\CAPETSHOW}[2]{%
    \def\CAPETargun{#1}%
    \ifx\CAPETargun\CAPETLIBELLEENONCES
      #1%            appel decoratif : rendu tel quel, hors numerotation
    \else
      \CAPETSHOWEXO{#1}{#2}%
    \fi
  }
  \newcommand{\CAPETSHOWEXO}[2]{%
    \stepcounter{CAPETexo}%
    \csgdef{CAPETcorrige\theCAPETexo}{#2}%
    \edef\CAPETtmpnum{\theCAPETexo}%
    #1%
    \par\vspace{2pt}%
    {\footnotesize\itshape\color{colExo}%
     \hyperref[capet:cor:\CAPETtmpnum]{Corrige de l'exercice \CAPETtmpnum\ $\rightarrow$ p.~\pageref{capet:cor:\CAPETtmpnum}}%
     \label{capet:eno:\CAPETtmpnum}}%
  }
  % \CAPETANNEXE : rejoue tous les corriges captures, dans l'ordre, avec intertitres
  \newcommand{\CAPETANNEXE}{%
    \setcounter{CAPETannexe}{0}%
    \loop\ifnum\value{CAPETannexe}<\value{CAPETexo}%
      \stepcounter{CAPETannexe}%
      \edef\CAPETcurnum{\theCAPETannexe}%
      \ifcsname CAPETnotionstart\CAPETcurnum\endcsname
        \par\medskip
        {\normalsize\bfseries\color{colExo}\csuse{CAPETnotionstart\CAPETcurnum}}%
        \par\nobreak\smallskip
      \fi
      \begin{tcolorbox}[exoBox, breakable,
        title={Corrige de l'exercice \CAPETcurnum\
               \hfill\normalfont\footnotesize
               \hyperref[capet:eno:\CAPETcurnum]{enonce p.~\pageref{capet:eno:\CAPETcurnum}}}]%
      \label{capet:cor:\CAPETcurnum}%
      \csuse{CAPETcorrige\CAPETcurnum}%
      \end{tcolorbox}%
    \repeat
  }
\else
  % --- Modes historiques, comportement strictement inchange ---
  \newcommand{\CAPETSHOW}[2]{\ifdefined\CAPETCORRIGE #2\else #1\fi}
  \newcommand{\CAPETANNEXE}{}
\fi

% ── Difficulté 1 à 3 (étoiles pleines, amssymb) ───────────────────────
\newcommand{\difficulte}[1]{%
  \ifcase#1\relax
  \or $\bigstar$\;niveau 1/3%
  \or $\bigstar\bigstar$\;niveau 2/3%
  \or $\bigstar\bigstar\bigstar$\;niveau 3/3%
  \fi
}

% ── Style encadré extrait de document technique (banque d'exercices) ──
\tcbset{
  docTech/.style={
    breakable,
    enhanced,
    colback=colGrayLight,
    colframe=colSeparator,
    fonttitle=\bfseries\small\itshape,
    coltitle=black,
    colbacktitle=colGrayLight,
    left=4pt, right=4pt, top=3pt, bottom=3pt,
    before skip=4pt, after skip=4pt,
  }
}

% ── Style tcolorbox de section principale ────────────────────────────
% Refonte 2026 : bandeau titre coloré arrondi + fond très pâle, sans cadre.
\tcbset{
  ficheSection/.style={
    breakable, enhanced jigsaw,
    frame hidden, boxrule=0pt, arc=1.8mm,
    colback=#1!4!white,
    colbacktitle=#1,
    coltitle=white,
    fonttitle=\sffamily\bfseries,
    left=7pt, right=7pt, top=5pt, bottom=5pt,
    toptitle=3pt, bottomtitle=3pt,
    before skip=10pt, after skip=8pt,
  }
}

% ── Style encadré question type ──────────────────────────────────────
% Carte blanche, bandeau sombre, liseré très clair.
\tcbset{
  questiontype/.style={
    breakable, enhanced jigsaw,
    colback=white,
    frame hidden, boxrule=0pt, arc=1.8mm,
    colbacktitle=colElec!78!black,
    coltitle=white,
    fonttitle=\sffamily\bfseries\small,
    left=7pt, right=7pt, top=4pt, bottom=4pt,
    toptitle=2.5pt, bottomtitle=2.5pt,
    before skip=8pt, after skip=8pt,
    borderline={0.6pt}{0pt}{colElec!12!white},
  }
}

% ── Style encadré jury ───────────────────────────────────────────────
% Bandeau ambre, fond crème, sans cadre.
\tcbset{
  jury/.style={
    breakable, enhanced jigsaw,
    colback=colJuryBg,
    frame hidden, boxrule=0pt, arc=1.8mm,
    colbacktitle=colJuryFrame,
    coltitle=white,
    fonttitle=\sffamily\bfseries,
    left=7pt, right=7pt, top=4pt, bottom=4pt,
    toptitle=3pt, bottomtitle=3pt,
    before skip=10pt, after skip=8pt,
  }
}

% ── En-tête / pied de page communs (refonte 2026) ─────────────────────
% #1 code fiche, #2 titre court, #3 couleur domaine
\newcommand{\ficheHeaders}[3]{%
  \fancyhead[L]{\small\sffamily{\bfseries\textcolor{#3}{#1}}\hspace{0.7em}\textcolor{colSlate}{#2}}%
  \fancyhead[R]{\small\sffamily\textcolor{colSlate}{CAPET SII IE\, --- \,Session 2026}}%
  \fancyfoot[C]{\small\sffamily\textcolor{colSlate}{\thepage\,/\,\pageref{LastPage}}}%
  \fancyhead[C]{}\fancyfoot[L]{}\fancyfoot[R]{}%
}

% ── Bloc en-tête de fiche (BLOC 0, refonte 2026) ──────────────────────
% #1 code  #2 titre  #3 couleur domaine  #4 domaine  #5 ligne présence  #6 sous-thèmes
\newcommand{\ficheEntete}[6]{%
  \begin{tcolorbox}[enhanced, breakable,
    frame hidden, boxrule=0pt, arc=2.5mm,
    interior style={left color=#3, right color=#3!75!black},
    left=9pt, right=9pt, top=7pt, bottom=7pt,
    before skip=0pt, after skip=6pt]
    {\color{white}\sffamily\bfseries\LARGE #1\quad #2}\\[4pt]
    {\color{#3!22!white}\sffamily\small \MakeUppercase{#4}\quad\textbar\quad #5}\\[1pt]
    {\color{#3!35!white}\sffamily\footnotesize Sous-thèmes : #6}
  \end{tcolorbox}%
}

% ── Mise en page ─────────────────────────────────────────────────────
\pagestyle{fancy}
\fancyhf{}
\renewcommand{\headrulewidth}{0.5pt}
\renewcommand{\footrulewidth}{0pt}
\AtBeginDocument{%
  \renewcommand{\headrule}{{\color{colRule}\hrule width\headwidth height 0.5pt}}%
}

\setlist[itemize]{leftmargin=1.2em, itemsep=1pt, topsep=2pt, parsep=0pt}
\setlist[enumerate]{leftmargin=1.2em, itemsep=3pt, topsep=2pt, parsep=0pt}

\setlength{\parskip}{2pt}
\setlength{\columnsep}{10pt}

\def\CAPETBEGINFILE{\begin{document}}
\def\CAPETENDFILE{\end{document}}
\else
\def\CAPETBEGINFILE{}
\def\CAPETENDFILE{}
\fi

\ficheHeaders{T20}{Filtres RC/RL}{colElec}

\CAPETBEGINFILE

% -----------------------------------------------------------------------
%  BLOC EN-TETE
% -----------------------------------------------------------------------

\ficheEntete{T20}{Filtres RC/RL - 1\textsuperscript{er} et 2\textsuperscript{e} ordre}{colElec}%
  {électronique - Traitement analogique du signal}%
  {\textbf{Présence :} \textbf{5/13} sujets (2013-2026) - thème \emph{régulier}}%
  {filtre passe-bas RC - filtre RL passe-bas - diagramme de Bode - gabarit fréquentiel - filtre du 2\textsuperscript{e} ordre (cascade, RLC) - filtre anti-repliement - Shannon}

% -----------------------------------------------------------------------
%  1. RAPPELS DE COURS
% -----------------------------------------------------------------------

\begin{tcolorbox}[
  ficheSection=colElec,
  title={1. Rappels de cours},
]

% -- 1.1 Filtre passe-bas RC du 1er ordre ------------------------------
\textbf{1.1 Filtre passe-bas RC - forme canonique}

\smallskip
Schéma : $R$ en série, $C$ en parallèle avec la sortie.
\[
  H(j\omega) = \frac{1}{1 + j\,\dfrac{\omega}{\omega_0}}
  \qquad \text{avec} \qquad \omega_0 = \frac{1}{RC}
\]
\begin{center}
\renewcommand{\arraystretch}{1.3}
\begin{tabular}{lll}
\toprule
\textbf{Grandeur} & \textbf{Expression} & \textbf{ $\omega = \omega_0$} \\
\midrule
Module $|H|$ & $\dfrac{1}{\sqrt{1+(\omega/\omega_0)^2}}$ & $\dfrac{1}{\sqrt{2}} \approx 0{,}707$ \\[6pt]
Gain en dB & $-20\log\sqrt{1+(\omega/\omega_0)^2}$ & $-3\,\mathrm{dB}$ \\[4pt]
Phase $\varphi$ & $-\arctan(\omega/\omega_0)$ & $-45^\circ$ \\[4pt]
Temps de réponse à 5\,\% & \multicolumn{2}{l}{$t_{r5\,\%} \approx 3\tau = 3\,RC$} \\
\bottomrule
\end{tabular}
\end{center}

\medskip\hrule\vspace{4pt}

% -- 1.2 Diagramme de Bode asymptotique --------------------------------
\textbf{1.2 Diagramme de Bode - règles asymptotiques}

\smallskip
\begin{center}
\renewcommand{\arraystretch}{1.3}
\begin{tabular}{lll}
\toprule
\textbf{Zone} & \textbf{Gain asymptotique} & \textbf{Phase asymptotique} \\
\midrule
$\omega \ll \omega_0$ & $0\,\mathrm{dB}$ (plateau) & $0^\circ$ \\
$\omega = \omega_0$ & $-3\,\mathrm{dB}$ (réel) & $-45^\circ$ (réel) \\
$\omega \gg \omega_0$ & pente $-20\,\mathrm{dB/decade}$ & $-90^\circ$ \\
\bottomrule
\end{tabular}
\end{center}

\smallskip
\begin{center}
\begin{tikzpicture}[scale=1, font=\footnotesize]
  \node[below] at (2.6,-2.15) {\bfseries Gain};
  \draw[->] (0,0) -- (5.4,0) node[right] {$\log\omega$};
  \draw[->] (0,-1.8) -- (0,1.3) node[above] {$G_{dB}$};
  \draw[thick, colElec] (0,0.85) -- (2,0.85);
  \draw[thick, colElec] (2,0.85) -- (5,-1.4);
  \draw[dashed] (2,0.85) -- (2,-1.8);
  \node[below] at (2,-1.8) {$\omega_0$};
  \node[above] at (0.9,0.85) {$0\,\mathrm{dB}$};
  \node[colElec] at (4.1,-0.55) {\rotatebox{-27}{$-20\,\mathrm{dB/dec}$}};

  \begin{scope}[xshift=7.4cm]
    \node[below] at (2.5,-2.15) {\bfseries Phase};
    \draw[->] (0,0.3) -- (5.4,0.3) node[right] {$\log\omega$};
    \draw[->] (0,-1.8) -- (0,0.6) node[above] {$\varphi$};
    \draw[thick, colElec] (0,0.3) .. controls (1.6,0.3) and (2.4,-1.35) .. (4,-1.35);
    \draw[thick, colElec] (4,-1.35) -- (5,-1.35);
    \draw[dashed] (2,0.3) -- (2,-1.8);
    \node[below] at (2,-1.8) {$\omega_0$};
    \node[left] at (-0.1,-0.68) {$-45^\circ$};
    \node[left] at (-0.1,-1.35) {$-90^\circ$};
  \end{scope}
\end{tikzpicture}
\end{center}

\medskip\hrule\vspace{4pt}

% -- 1.3 Filtre RL passe-bas -------------------------------------------
\textbf{1.3 Filtre RL passe-bas}

\smallskip
$R$ en parallèle avec $L$ : même forme que RC avec $\tau = L/R$.
\[
  H(j\omega) = \frac{1}{1 + j\,\dfrac{\omega}{\omega_0}}
  \qquad \omega_0 = \frac{R}{L}
\]
Les enroulements d'une machine asynchrone se comportent comme un filtre RL passe-bas naturel qui atténue les harmoniques de la MLI.

\medskip\hrule\vspace{4pt}

% -- 1.4 Filtre du 2e ordre --------------------------------------------
\textbf{1.4 Filtre passe-bas du 2\textsuperscript{e} ordre - forme canonique}

\[
  H(j\omega) = \frac{H_0}{1 + 2m\,\dfrac{j\omega}{\omega_0} + \left(\dfrac{j\omega}{\omega_0}\right)^2}
\]

\begin{center}
\renewcommand{\arraystretch}{1.3}
\begin{tabular}{lll}
\toprule
\textbf{Paramètre} & \textbf{Signification} & \textbf{Effet} \\
\midrule
$\omega_0$ & Pulsation propre (rad/s) & Position de la coupure \\
$m$ & Facteur d'amortissement & Forme de la réponse \\
$Q = 1/(2m)$ & Facteur de qualité & Sélectivité \\
\midrule
$m > 1$ & Sur-amorti & Pas d'oscillations \\
$m = 1$ & Critique & Retour le plus rapide sans dépassement \\
$m < 1$ & Sous-amorti & Résonance, pic à $\omega_r = \omega_0\sqrt{1-2m^2}$ \\
\bottomrule
\end{tabular}
\end{center}

\medskip\hrule\vspace{4pt}

% -- 1.5 Cascade de 2 filtres 1er ordre --------------------------------
\textbf{1.5 Cascade de deux filtres RC identiques}

\[
  H(j\omega) = \left(\frac{1}{1 + j\omega/\omega_0}\right)^2
\]
Pente asymptotique $-40\,\mathrm{dB/decade}$ ; amortissement $m = 1$ (pas d'oscillations ni de résonance).

\medskip\hrule\vspace{4pt}

% -- 1.6 Filtre anti-repliement ----------------------------------------
\textbf{1.6 Filtre anti-repliement (Shannon)}

Placé avant le CAN pour éliminer les composantes à $f > f_e/2$ :
\[
  f_e \geq 2\,f_{max} \qquad \Rightarrow \qquad f_c \leq \frac{f_e}{2}
\]

\end{tcolorbox}

% -----------------------------------------------------------------------
%  2. FORMULAIRE SYNTHETIQUE
% -----------------------------------------------------------------------

\begin{tcolorbox}[
  ficheSection=colElec,
  title={2. Formulaire synthétique},
]
\small
\begin{multicols}{2}
\renewcommand{\arraystretch}{1.45}
\begin{tabular}{@{}p{4.8cm}ll@{}}
\toprule
\textbf{Formule} & \textbf{Usage} & \textbf{U.} \\
\midrule
\multicolumn{3}{@{}l}{\footnotesize\textit{Filtre 1\textsuperscript{er} ordre RC}} \\
$\omega_0 = 1/(RC)$ & Coupure & rad/s \\
$|H(\omega_0)| = 1/\sqrt{2}$ & Module à $\omega_0$ & - \\
$G_{dB}(\omega_0) = -3\,\mathrm{dB}$ & Gain & dB \\
$\varphi(\omega_0) = -45^\circ$ & Phase & - \\
$t_{r5\%} \approx 3\,RC$ & Transitoire & s \\
\midrule
\multicolumn{3}{@{}l}{\footnotesize\textit{Filtre RL}} \\
$\omega_0 = R/L$ & Coupure & rad/s \\
$\tau = L/R$ & Cste de temps & s \\
\bottomrule
\end{tabular}

\columnbreak

\renewcommand{\arraystretch}{1.45}
\begin{tabular}{@{}p{4.8cm}ll@{}}
\toprule
\textbf{Formule} & \textbf{Usage} & \textbf{U.} \\
\midrule
\multicolumn{3}{@{}l}{\footnotesize\textit{Filtre 2\textsuperscript{e} ordre}} \\
$H = H_0 / (1 + 2m\,(j\omega/\omega_0) + (j\omega/\omega_0)^2)$ & Forme canon. & - \\
$Q = 1/(2m)$ & Qualité & - \\
$\omega_r = \omega_0\sqrt{1-2m^2}$ & Résonance ($m<1/\sqrt{2}$) & rad/s \\
Pente HF : $-40\,\mathrm{dB/dec}$ & Bode & - \\
\midrule
\multicolumn{3}{@{}l}{\footnotesize\textit{Shannon / anti-repliement}} \\
$f_e \geq 2\,f_{max}$ & Shannon & Hz \\
$f_c \leq f_e/2$ & Filtre AR & Hz \\
\midrule
\multicolumn{3}{@{}l}{\footnotesize\textit{Gabarit}} \\
Atténuation $>100$ ($>40\,\mathrm{dB}$) $\Rightarrow$ 2\textsuperscript{e} ordre & & \\
\bottomrule
\end{tabular}
\end{multicols}
\normalsize
\end{tcolorbox}

% -----------------------------------------------------------------------
%  3. QUESTIONS TYPES
% -----------------------------------------------------------------------

\begin{tcolorbox}[
  ficheSection=colElec,
  title={3. Questions types - extraites des sujets réels},
]

% -- QT1 -------------------------------------------------------------

\begin{tcolorbox}[questiontype,
  title={QT1 - Filtre RC - fonction de transfert et diagramme de Bode \hfill\normalfont\footnotesize E1 2026 (B3-B4)}
]
\textbf{Énoncé.}
Filtre passe-bas RC avec $R = \SI{1}{\kilo\ohm}$ et $C = \SI{100}{\nano\farad}$. Montrer que $H_{F1}(j\omega) = 1/(1+j\omega/\omega_0)$. Exprimer $\omega_0$. Tracer le diagramme de Bode asymptotique.

\tcbline

\textbf{Corrigé.}

\textbf{Étape 1 — Diviseur de tension.} ($Z_C = 1/(jC\omega)$) :
\[
  H = \frac{Z_C}{R + Z_C} = \frac{1/(jC\omega)}{R + 1/(jC\omega)} = \frac{1}{1 + jRC\omega}
  \qquad \Rightarrow \qquad \omega_0 = \frac{1}{RC} = \frac{1}{10^3 \times 10^{-7}} = \SI{10000}{\radian\per\second}
\]

\textbf{Étape 2 — Bode asymptotique.} Plateau $0\,\mathrm{dB}$ pour $\omega \ll \omega_0$ ; pente $-20\,\mathrm{dB/decade}$ pour $\omega \gg \omega_0$ ; point réel à $\omega_0$ : $-3\,\mathrm{dB}$, $-45^\circ$.

\textbf{Conclusion.} La coupure à $\omega_0 = 10\,000\,\mathrm{rad/s}$ ($f_0 \approx 1592\,\mathrm{Hz}$) situe ce filtre pour couper les harmoniques au-delà du kHz tout en laissant passer le continu et le signal utile en basse fréquence.
\end{tcolorbox}

\vspace{4pt}

% -- QT2 -------------------------------------------------------------

\begin{tcolorbox}[questiontype,
  title={QT2 - Lecture graphique sur le Bode - calcul de $v_3(t)$ \hfill\normalfont\footnotesize E1 2026 (B5-B6)}
]
\textbf{Énoncé.}
Signal d'entrée $v_2(t) = 2 + 0{,}5\sin(2\pi \cdot 900 \cdot t)$. Avec $f_0 = \SI{1592}{\hertz}$ ($\omega_0 = 10^4\,\mathrm{rad/s}$), déterminer le signal de sortie $v_3(t)$ et estimer le temps de réponse à 5\,\%.

\tcbline

\textbf{Corrigé.}

\textbf{Étape 1 — Composante DC.} $\omega = 0$ : $|H(0)| = 1$ - $V_{DC} = 2\,\mathrm{V}$.

\textbf{Étape 2 — Composante AC.} $f_2 = 900\,\mathrm{Hz}$ - $\omega_2 = 2\pi \times 900 \approx 5655\,\mathrm{rad/s}$ :
\[
  |H(\omega_2)| = \frac{1}{\sqrt{1+(5655/10000)^2}} = \frac{1}{\sqrt{1+0{,}32}} \approx 0{,}870
  \qquad
  V_{AC,out} = 0{,}5 \times 0{,}870 \approx 0{,}435\,\mathrm{V}
\]
\[
  v_3(t) \approx 2 + 0{,}435\sin(2\pi \times 900\,t + \varphi) \,\mathrm{V}
\]

\textbf{Étape 3 — Temps de réponse.} $t_{r5\%} \approx 3\,RC = 3 \times 10^3 \times 10^{-7} = 0{,}3\,\mathrm{ms}$.

\textbf{Conclusion.} Le cahier des charges de rapidité ($< 50\,\mathrm{ms}$) est largement satisfait, mais l'atténuation à 900\,Hz reste insuffisante (filtre du 1\textsuperscript{er} ordre) : un filtre du 2\textsuperscript{e} ordre est nécessaire pour mieux rejeter cette composante.
\end{tcolorbox}

\vspace{4pt}

% -- QT3 -------------------------------------------------------------

\begin{tcolorbox}[questiontype,
  title={QT3 - Filtre du 2\textsuperscript{e} ordre - choix des composants et réponse indicielle \hfill\normalfont\footnotesize E1 2026 (B9)}
]
\textbf{Énoncé.}
On souhaite atténuer $> 40\,\mathrm{dB}$ à $\omega_2 = 5655\,\mathrm{rad/s}$ avec un temps de réponse à 5\,\% $< 50\,\mathrm{ms}$. Proposer $R$ et $C$ pour un filtre 2\textsuperscript{e} ordre (cascade de deux RC). Justifier si la réponse indicielle présente des oscillations.

\tcbline

\textbf{Corrigé.}

\textbf{Étape 1 — Condition de gain.} $|H(\omega_2)|^2 = [1/(1+(\omega_2/\omega_0)^2)] \leq 1/100$ - $\omega_0 \leq \omega_2/\sqrt{99} \approx 568\,\mathrm{rad/s}$ - $f_0 \leq 90\,\mathrm{Hz}$.

\textbf{Étape 2 — Condition de rapidité.} $t_{r5\%} \approx 3RC \leq 50\,\mathrm{ms}$ - $RC \leq 16{,}7\,\mathrm{ms}$.

\textbf{Étape 3 — Choix des composants.} $RC = 1/(568) \approx 1{,}76\,\mathrm{ms}$ - par ex. $R = \SI{17{,}6}{\kilo\ohm}$, $C = \SI{100}{\nano\farad}$.

\textbf{Conclusion.} La cascade de deux filtres RC identiques donne $m = 1$ (régime critique) : \textbf{pas d'oscillations} sur la réponse indicielle, ce qui satisfait à la fois le gabarit fréquentiel et la contrainte de stabilité du signal de sortie.
\end{tcolorbox}

\end{tcolorbox}

% -----------------------------------------------------------------------
%  4. ATTENTES DU JURY
% -----------------------------------------------------------------------

\begin{tcolorbox}[
  jury,
  title={4. Attentes du jury - Erreurs récurrentes et points de vigilance},
]

\begin{enumerate}[label=\textbf{\arabic*.}, leftmargin=*, itemsep=5pt]

  \item \textbf{Maîtriser la forme canonique} $H = 1/(1+j\omega/\omega_0)$ et l'identifier dans tout schéma RC ou RL - ne pas recalculer depuis les équations de circuit si la forme est déjà fournie.

  \item \textbf{Ne pas confondre $\omega_0$ et $f_0$} : $\omega_0 = 2\pi f_0$ en rad/s, $f_0$ en Hz. Donner les unités explicitement dans chaque résultat.

  \item \textbf{Gain à $\omega_0$ vaut $-3\,\mathrm{dB}$, pas 0\,dB.} Erreur fréquente : les candidats écrivent $|H(\omega_0)| = 0$ au lieu de $1/\sqrt{2}$.

  \item \textbf{Justifier l'ordre du filtre depuis le gabarit} : pente disponible par ordre = $-20\,\mathrm{dB/decade}$. Pour atténuer $>40\,\mathrm{dB}$ - 2\textsuperscript{e} ordre minimum. Ne pas choisir l'ordre par intuition.

  \item \textbf{Filtre anti-repliement} : systématiquement mentionné dès qu'un CAN apparaît dans le schéma. Rappeler le théorème de Shannon, puis déduire $f_c \leq f_e/2$.

  \item \textbf{Cascade vs RLC} : deux filtres RC en cascade donnent $m = 1$ (pas de résonance), contrairement à un circuit RLC série ou parallèle.

\end{enumerate}

\medskip
\begin{tcolorbox}[
  colback=colJuryFrame!12!colJury,
  colframe=colJuryFrame,
  left=4pt, right=4pt, top=2pt, bottom=2pt,
  fontupper=\small\itshape
]
Rapport jury 2023 : «~La quatrième partie portait sur l'asservissement, trop peu de candidats l'ont traitée de manière satisfaisante. [\ldots] Le jury invite les candidats à revoir les concepts de base de ce domaine comme les notions de boucle ouverte, de correcteur ou encore de marge de phase.~»\\
Corrigé officiel 2024 (question B13), et non commentaire de jury : «~la suppression des harmoniques en haute fréquence est faisable par l'effet inductif des enroulements de la machine.~»
\end{tcolorbox}

\end{tcolorbox}

% -----------------------------------------------------------------------
%  PIED - PRESENCE DANS LES SUJETS
% -----------------------------------------------------------------------

\begin{tcolorbox}[
  ficheSection=colPresence,
  title={Présence dans les sujets analysés (2013-2026)},
]
\begin{center}
\footnotesize\renewcommand{\arraystretch}{1.15}
\begin{tabular}{lccccccccccccc}
\toprule
\textbf{Année} & 13 & 14 & 15 & 16 & 17 & 18 & 19 & 20 & 21 & 22 & 23 & 24 & 25 \\
\midrule
E1 (commune)   & - & - & - & - & - & - & - & - & - & - & $\checkmark$ & - & $\checkmark$ \\
E2 discipl. IE & - & - & - & - & - & $\checkmark$ & - & $\checkmark$ & - & - & - & $\checkmark$ & - \\
E2 appl. comm. & - & - & - & - & - & - & - & - & - & - & - & - & - \\
\bottomrule
\end{tabular}
\end{center}
\smallskip
\footnotesize\textbf{Sous-thèmes fréquents :}
filtre passe-bas RC ($\omega_0$, Bode) - diagramme de Bode asymptotique + réel - gabarit fréquentiel - filtre anti-repliement + Shannon - cascade de 2 filtres RC

\medskip
\textbf{Ordres de grandeur (repères rapides).}
\begin{itemize}[leftmargin=1.3em]
  \item Réseau BT FR : tension phase-neutre typique \SI{230}{\volt}, tension phase-phase typique \SI{400}{\volt}, fréquence \SI{50}{\hertz}.
  \item Courants d'étude fréquents en sujet : de quelques \si{\ampere} à quelques dizaines de \si{\ampere}.
  \item En conversion d'énergie, un rendement attendu reste généralement entre 0{,}8 et 0{,}98 selon la technologie.
\end{itemize}

\medskip
\textbf{Synthèse express (à refaire sans support).}
\begin{enumerate}[leftmargin=1.4em]
  \item Poser les hypothèses du modèle et les conventions de signe.
  \item Écrire 2 à 4 relations clés avec unités.
  \item Réaliser un calcul numérique complet et vérifier l'ordre de grandeur.
  \item Justifier le résultat par un argument technique de niveau CAPET.
\end{enumerate}

\end{tcolorbox}

\CAPETENDFILE
